% ══════════════════════════════════════════════════════════════════════
% Chapter 1: Hermeneutic Foundations
% ══════════════════════════════════════════════════════════════════════

\chapter{Hermeneutic Foundations: Repertoires, Interpretation, and the Hermeneutic Bound}
\label{ch:hermeneutic}

\epigraph{``The limits of my language mean the limits of my world.''}{Ludwig Wittgenstein, \textit{Tractatus Logico-Philosophicus} (5.6)}

% ── Introduction ──────────────────────────────────────────────────────

Every human mind is a repertoire: a stock of ideas accumulated through
enculturation.  When a signal arrives---a gesture, an utterance, a ritual
act---the receiver does not passively absorb it.  They \emph{interpret}
it by composing it with their existing ideas.  This chapter builds the
foundational structures on which the entire sequel rests.

In the language of Bourdieu~\cite{bourdieu1984}, a repertoire is the
cognitive component of \emph{habitus}---the durable, transposable
dispositions through which agents perceive and act in the social world.
In the language of Saussure~\cite{saussure1916}, each element of the
repertoire is a signified, and interpretation is the operation that maps
a signifier to a new signified by composing it with an existing one.  In
the language of Peirce~\cite{peirce1931}, interpretation is the
generation of an \emph{interpretant} from the encounter of a sign with a
prior sign.

The payoff in a signalling game is \emph{never} about the signal
itself---it is always about the interpretations the signal produces in
other minds.  This insight, already present in Gadamer's
\emph{Truth and Method}~\cite{gadamer1960}, receives its first fully
formal treatment here.

\medskip
\noindent\textbf{Chapter overview.}  We introduce four core concepts:
\emph{repertoire} (\S\ref{sec:ch1:core}), \emph{interpretation} and the
\emph{hermeneutic bound} (\S\ref{sec:ch1:hermeneutic-bound}), \emph{silence}
(\S\ref{sec:ch1:silence}), and \emph{cultural consensus}
(\S\ref{sec:ch1:consensus}).
We close with the \emph{telephone game} (\S\ref{sec:ch1:telephone}) and
structural invariants (\S\ref{sec:ch1:invariants}).  All sixteen theorems
are machine-verified in Lean~4 with zero \texttt{sorry} axioms.

% ── §1. Core Definitions ─────────────────────────────────────────────
\section{Core Definitions}
\label{sec:ch1:core}

We work within an ideatic space $(\IS, \compose, \void, \depth, \res)$
as established in the IDT Foundations.

\begin{definition}[Repertoire]
\label{def:repertoire}
A \emph{repertoire} is a finite list of ideas:
\[
  \Rep = [r_1, r_2, \ldots, r_n] \in \IS^{*}.
\]
In Bourdieu's terms, this is the \emph{habitus} made formal---the
totality of schemas through which an agent perceives and interprets the
social world.
\end{definition}

\begin{definition}[Interpretation]
\label{def:interpretation}
When a receiver with repertoire $\Rep = [r_1, \ldots, r_n]$ encounters
signal~$s$, they produce the list of interpretations
\[
  \operatorname{interpret}(\Rep, s) 
  = [r_1 \compose s,\; r_2 \compose s,\; \ldots,\; r_n \compose s].
\]
This is Gadamer's \emph{Horizontverschmelzung} (fusion of horizons):
each background idea merges with the incoming signal.
\end{definition}

In Lean~4, the definition is:

\begin{lstlisting}[caption={Interpretation in Lean~4},label=lst:interpret]
def interpret (rep : Repertoire I) (signal : I) : List I :=
  rep.map (fun r => IdeaticSpace.compose r signal)
\end{lstlisting}

\begin{definition}[Maximum Interpretation Depth]
\label{def:maxinterpdepth}
The \emph{maximum interpretation depth} of a list of ideas is
\[
  \operatorname{maxInterpDepth}([a_1, \ldots, a_n])
  = \max_{i} \depth(a_i).
\]
\end{definition}

\begin{definition}[Depth Sum]
\label{def:depthsum}
The \emph{depth sum} of a list of ideas is
$\operatorname{depthSum}([a_1, \ldots, a_n]) = \sum_i \depth(a_i)$.
\end{definition}

% ── §2. The Hermeneutic Bound ────────────────────────────────────────
\section{The Hermeneutic Bound}
\label{sec:ch1:hermeneutic-bound}

``You can only understand what you are prepared to understand.''
The following two theorems make this aphorism precise.

\begin{theorem}[Total Hermeneutic Bound]
\label{thm:total-hermeneutic-bound}
For any repertoire $\Rep$ and signal~$s$,
\[
  \operatorname{depthSum}(\operatorname{interpret}(\Rep, s))
  \;\leq\;
  \operatorname{depthSum}(\Rep) + |\Rep| \cdot \depth(s).
\]
\end{theorem}

\begin{proof}[Proof sketch]
By induction on $|\Rep|$.  The base case ($\Rep = []$) is trivial.
For $\Rep = r :: \Rep'$, we use the depth-composition bound
$\depth(r \compose s) \leq \depth(r) + \depth(s)$ and the
inductive hypothesis on~$\Rep'$.
The full proof is verified in \texttt{IDT\_Signal\_Ch01.lean},
theorem \texttt{total\_hermeneutic\_bound}.
\end{proof}

\begin{remark}[Mass Media]
This theorem explains why mass media (one signal, many receivers) has
bounded total interpretive impact per receiver.  A richer repertoire
absorbs more, but the signal's contribution scales only linearly with
audience size.
\end{remark}

\begin{theorem}[MaxDepth Interpretation Bound]
\label{thm:max-interp-depth}
For any repertoire $\Rep$ and signal~$s$,
\[
  \operatorname{maxInterpDepth}(\operatorname{interpret}(\Rep, s))
  \;\leq\;
  \operatorname{maxInterpDepth}(\Rep) + \depth(s).
\]
\end{theorem}

\begin{proof}[Proof sketch]
By induction on $|\Rep|$.  At each step, we bound
$\depth(r \compose s) \leq \depth(r) + \depth(s)$ and propagate
through the maximum.  See theorem \texttt{max\_interp\_depth\_bound}.
\end{proof}

\begin{remark}
Even the most sophisticated listener adds only $\depth(s)$ to their
best existing understanding.  No single signal can vault a receiver
beyond their current horizon by more than the signal's intrinsic
complexity.
\end{remark}

% ── §3. Silence and the Unsaid ───────────────────────────────────────
\section{Silence and the Unsaid}
\label{sec:ch1:silence}

\begin{theorem}[Silence is Transparent]
\label{thm:silence-transparent}
The void signal leaves the mind unchanged:
\[
  \operatorname{interpret}(\Rep, \void) = \Rep.
\]
\end{theorem}

\begin{proof}[Proof sketch]
By induction on $\Rep$.  For each $r \in \Rep$, the void-right identity
$r \compose \void = r$ yields the result immediately.
See \texttt{silence\_is\_transparent}.
\end{proof}

\begin{remark}[Ritual Silence]
Quaker meetings, Buddhist \emph{noble silence}, the minute of silence
for the dead---all exploit this principle.  By saying nothing, the
receiver's existing framework remains undisturbed.  The void signal is
the formal foundation of taboo: what must not be said leaves the mind in
its current state.
\end{remark}

\begin{theorem}[The Blank Receiver]
\label{thm:blank-receiver}
A receiver whose only idea is $\void$ interprets any signal as the
signal itself:
\[
  \operatorname{interpret}([\void], s) = [s].
\]
\end{theorem}

\begin{proof}[Proof sketch]
Direct computation: $\void \compose s = s$ by the void-left identity.
See \texttt{blank\_receiver}.
\end{proof}

\begin{remark}[The Objective Observer]
The impossible ideal of the ``objective observer'' requires total
cultural amnesia: a repertoire consisting solely of~$\void$.  Only then
does interpretation reduce to transparent reception.
\end{remark}

\begin{theorem}[Negative Capability / Keats's Theorem]
\label{thm:negative-capability}
If $\void \in \Rep$, then $s \in \operatorname{interpret}(\Rep, s)$.
\end{theorem}

\begin{proof}[Proof sketch]
Since $\void \in \Rep$, the map sends $\void \mapsto \void \compose s = s$,
so $s$ appears in the output.
See \texttt{negative\_capability}.
\end{proof}

\begin{remark}
Keats described \emph{negative capability} as the capacity for ``being
in uncertainties, Mysteries, doubts, without any irritable reaching after
fact \& reason'' \cite{keats1817}.  In our framework, a mind retaining a
slot of pure receptivity ($\void$ in the repertoire) captures the signal
in its original form alongside culturally mediated versions.
\end{remark}

% ── §4. Cultural Consensus ───────────────────────────────────────────
\section{Cultural Consensus}
\label{sec:ch1:consensus}

The next three theorems establish that resonance in backgrounds implies
resonance in interpretations---the structural foundation of cultural
consensus.

\begin{theorem}[Durkheimian Consensus]
\label{thm:durkheimian-consensus}
If $r_1 \res r_2$ and both receive signal~$s$, then
$r_1 \compose s \;\res\; r_2 \compose s$.
\end{theorem}

\begin{proof}[Proof sketch]
Direct application of the resonance-compose-right lemma from IDT
Foundations.  See \texttt{durkheimian\_consensus}.
\end{proof}

\begin{remark}
Two members of the same culture (shared myths, shared categories)
\emph{always} produce resonant interpretations of the same event.
Cultural consensus is a structural consequence of shared
background---Durkheim's \emph{conscience collective}~\cite{durkheim1893}.
\end{remark}

\begin{theorem}[L\'evi-Strauss's Mythical Transformation]
\label{thm:mythical-transformation}
For any receiver~$r$, if $s_1 \res s_2$ then
$r \compose s_1 \;\res\; r \compose s_2$.
\end{theorem}

\begin{proof}[Proof sketch]
By resonance-compose-left.  See \texttt{mythical\_transformation}.
\end{proof}

\begin{remark}
Ritual repetition-with-variation works because the participant's
invariant background ensures that structurally similar performances
produce resonant experiences.  L\'evi-Strauss's mythical
``transformations'' are perceived as ``the same story told differently''
precisely because the listener's background resonates with both
variants~\cite{levistrauss1962}.
\end{remark}

\begin{theorem}[Intercultural Solidarity / Turner's Communitas]
\label{thm:turners-communitas}
If $r_1 \res r_2$ and $s_1 \res s_2$, then
\[
  r_1 \compose s_1 \;\res\; r_2 \compose s_2.
\]
\end{theorem}

\begin{proof}[Proof sketch]
By the full resonance-compose lemma.  See \texttt{turners\_communitas}.
\end{proof}

\begin{remark}
Members of allied cultures witnessing similar rituals produce resonant
interpretations.  This is the formal basis of what Turner called
\emph{communitas}~\cite{turner1969}: the spontaneous solidarity that
arises when structurally similar participants undergo similar experiences.
\end{remark}

% ── §5. The Telephone Game ───────────────────────────────────────────
\section{The Telephone Game: Cultural Transmission Chains}
\label{sec:ch1:telephone}

\begin{theorem}[Bartlett's Serial Reproduction]
\label{thm:bartlett}
Two-stage interpretation through different minds equals single-stage
interpretation through the composed mind:
\[
  r_B \compose (r_A \compose s)
  \;=\;
  (r_B \compose r_A) \compose s.
\]
\end{theorem}

\begin{proof}[Proof sketch]
This is the associativity axiom of the ideatic space, applied to the
triple $(r_B, r_A, s)$.
See \texttt{bartletts\_serial\_reproduction}.
\end{proof}

\begin{remark}
Bartlett's (1932) serial reproduction experiments~\cite{bartlett1932}
showed that stories change predictably through chains of retelling.  This
theorem shows \emph{why}: the chain of interpreters is structurally
equivalent to a single composite interpreter.  Oral tradition scholars
can study the ``net effect'' of a transmission chain without tracking
each individual retelling.
\end{remark}

\begin{theorem}[Telephone Depth Bound]
\label{thm:telephone-depth}
The depth after two-stage interpretation satisfies
\[
  \depth(r_B \compose (r_A \compose s))
  \;\leq\;
  \depth(r_B) + \depth(r_A) + \depth(s).
\]
\end{theorem}

\begin{proof}[Proof sketch]
Rewrite using associativity, then apply depth-compose-le twice.
See \texttt{telephone\_depth}.
\end{proof}

\begin{theorem}[Resonant Transmission Lines]
\label{thm:resonant-transmission}
If $r_{A_1} \res r_{A_2}$ and $r_{B_1} \res r_{B_2}$, then
\[
  r_{B_1} \compose (r_{A_1} \compose s)
  \;\res\;
  r_{B_2} \compose (r_{A_2} \compose s).
\]
\end{theorem}

\begin{proof}[Proof sketch]
Apply associativity to both sides, then use resonance-compose twice.
See \texttt{resonant\_transmission\_lines}.
\end{proof}

\begin{remark}
Independent lines of ritual transmission (e.g., different griots in West
Africa, different monastic orders transmitting the same texts) produce
resonant outcomes if corresponding specialists resonate.  Ritual
orthodoxy can be maintained across geographically separated communities
without central authority.
\end{remark}

% ── §6. Structural Invariants ────────────────────────────────────────
\section{Structural Invariants of Interpretation}
\label{sec:ch1:invariants}

\begin{theorem}[Void is Unpersuasive]
\label{thm:void-unpersuasive}
$\depth(r \compose \void) = \depth(r)$.  Silence contributes zero
marginal depth.
\end{theorem}

\begin{proof}[Proof sketch]
By the depth-compose-void identity.  See \texttt{void\_is\_unpersuasive}.
\end{proof}

\begin{remark}[Game-theoretic significance]
The null message (babbling) contributes zero information.  This is
\emph{why} babbling equilibria always exist: silence costs nothing and
changes nothing.
\end{remark}

\begin{theorem}[Interpretation Conserves Count]
\label{thm:conserves-count}
$|\operatorname{interpret}(\Rep, s)| = |\Rep|$.
\end{theorem}

\begin{proof}[Proof sketch]
By the length-preserving property of \texttt{List.map}.
See \texttt{interpretation\_conserves\_count}.
\end{proof}

\begin{theorem}[Durkheim's Division of Labor]
\label{thm:division-of-labor}
$\operatorname{interpret}(\Rep_1 \mathbin{+\!+} \Rep_2, s)
 = \operatorname{interpret}(\Rep_1, s) \mathbin{+\!+}
   \operatorname{interpret}(\Rep_2, s)$.
\end{theorem}

\begin{proof}[Proof sketch]
By \texttt{List.map\_append}.  See \texttt{durkheims\_division\_of\_labor}.
\end{proof}

\begin{remark}
A polymath's interpretation equals the union of specialist
interpretations.  Interdisciplinary understanding is \emph{not}
synergistic---it is the disjoint union of specialist understandings.
This formalizes Durkheim's cognitive division of labor~\cite{durkheim1893}.
\end{remark}

\begin{theorem}[Silence Preserves DepthSum]
\label{thm:silence-depthsum}
$\operatorname{depthSum}(\operatorname{interpret}(\Rep, \void))
 = \operatorname{depthSum}(\Rep)$.
\end{theorem}

\begin{proof}[Proof sketch]
Follows immediately from Theorem~\ref{thm:silence-transparent}.
See \texttt{silence\_preserves\_depthSum}.
\end{proof}

\begin{theorem}[Three-Stage Telephone]
\label{thm:three-stage}
\[
  r_C \compose (r_B \compose (r_A \compose s))
  \;=\;
  ((r_C \compose r_B) \compose r_A) \compose s.
\]
\end{theorem}

\begin{proof}[Proof sketch]
Three applications of associativity.
See \texttt{three\_stage\_telephone}.
\end{proof}

\begin{remark}
Arbitrarily long chains of cultural transmission collapse into a single
composite interpreter.  This is why we can meaningfully speak of ``the
Western tradition'' as a single interpretive lens, despite it being the
product of thousands of individual acts of interpretation.
\end{remark}

% ── Summary ──────────────────────────────────────────────────────────
\section*{Chapter Summary}

This chapter established the foundational apparatus of the signalling-game
framework within IDT\@.  We introduced \emph{repertoires} as formal models
of culturally constituted minds, and \emph{interpretation} as
compositional fusion of horizons.  The \emph{Total Hermeneutic Bound}
(Theorem~\ref{thm:total-hermeneutic-bound}) shows that no signal can
inject more complexity into a mind than the signal's depth times the
repertoire's size.  \emph{Silence is Transparent}
(Theorem~\ref{thm:silence-transparent}) provides the formal foundation
for babbling equilibria.  The \emph{Durkheimian Consensus}
(Theorem~\ref{thm:durkheimian-consensus}) and \emph{Turner's Communitas}
(Theorem~\ref{thm:turners-communitas}) establish that resonant
backgrounds guarantee resonant interpretations, providing the formal
basis of cultural consensus.  The \emph{Telephone Game} theorems
(Theorems~\ref{thm:bartlett}--\ref{thm:resonant-transmission}) show that
sequential interpretation chains collapse into single composite
interpreters via associativity.

All sixteen theorems are verified in
\texttt{FormalAnthropology/IDT\_Signal\_Ch01.lean}.
